\begin{tikzpicture}[
  x=0.72cm,
  y=0.72cm,
  line cap=round,
  line join=round,
  positive/.style={circle,fill=black,inner sep=1.65pt},
  negative/.style={circle,draw=black,fill=white,inner sep=1.15pt,
    line width=0.45pt},
  reference/.style={draw=black!55,line width=0.55pt},
  interface/.style={draw=black,line width=1.25pt},
  sector/.style={-{Latex[length=2.2mm,width=1.25mm]},line width=0.65pt},
  ring/.style={draw=black!55,line width=0.65pt},
  slip/.style={fill=black},
  panel/.style={font=\small\bfseries},
  note/.style={font=\scriptsize,align=center}
]
  % Panel (a): a reference gap and the elementary G6 phase slip.
  \node[panel,anchor=west] at (-0.35,3.55) {(a)};
  \node[note,anchor=east] at (-0.35,2.75) {reference};
  \foreach \i in {0,...,4} {
    \coordinate (r\i) at (\i,2.75);
  }
  \draw[reference] (r0)--(r4);
  \node[positive] at (r0) {};
  \node[positive] at (r4) {};
  \foreach \i in {1,2,3} {
    \node[negative] at (r\i) {};
  }
  \draw[<->,line width=0.45pt]
    ($(r0)+(0,-0.36)$)--node[note,fill=white,inner sep=1pt]
    {$g=4,\ q=0$}($(r4)+(0,-0.36)$);

  \node[note,anchor=east] at (-0.35,1.10) {G6};
  \foreach \i in {0,...,6} {
    \coordinate (g\i) at ({2*\i/3},1.10);
  }
  \draw[interface] (g0)--(g6);
  \node[positive] at (g0) {};
  \node[positive] at (g6) {};
  \foreach \i in {1,...,5} {
    \node[negative] at (g\i) {};
  }
  \draw[<->,line width=0.45pt]
    ($(g0)+(0,-0.36)$)--node[note,fill=white,inner sep=1pt]
    {$g=6,\ q=2$}($(g6)+(0,-0.36)$);
  \draw[sector] (0.75,-0.18)--node[note,above]
    {$B_s\longmapsto B_{s+2}$} (3.25,-0.18);

  % Panel (b): legal residue rings. Thick ticks are well-separated G6 cores.
  \node[panel,anchor=west] at (5.10,3.55) {(b)};
  \foreach \cx/\rslips/\residue/\sectorchange in {
    6.10/1/2/2,
    9.10/2/4/0,
    12.10/3/6/2
  } {
    \draw[ring] (\cx,2.25) circle (0.94);
    \foreach \j in {1,...,\rslips} {
      \pgfmathsetmacro{\ang}{90+360*(\j-1)/\rslips}
      \pgfmathsetmacro{\xa}{\cx+0.78*cos(\ang)}
      \pgfmathsetmacro{\ya}{2.25+0.78*sin(\ang)}
      \pgfmathsetmacro{\xb}{\cx+1.10*cos(\ang)}
      \pgfmathsetmacro{\yb}{2.25+1.10*sin(\ang)}
      \draw[interface] (\xa,\ya)--(\xb,\yb);
      \fill[slip] (\xb,\yb) circle (1.45pt);
    }
    \node[note] at (\cx,2.25) {$\rslips\times\mathrm{G6}$};
    \node[note] at (\cx,0.74)
      {\(n\bmod8=\residue\)\\
       \(q_{\rm tot}=\residue,\ \Delta s=\sectorchange\)};
  }
\end{tikzpicture}
